\section{Neural Dynamics Model}

\subsection{History Encoder}

Let the history window be
\begin{equation}
\mathbf{X}^{(L)}_k \in \mathbb{R}^{L \times D_{\mathrm{in}}}.
\end{equation}

The current baseline model family is \texttt{S1Predictor}, which uses a sequence encoder to summarize the control and proxy-state history:
\begin{equation}
\mathbf{h}_k = \Phi_{\theta}\!\left(\mathbf{X}^{(L)}_k\right).
\end{equation}

In the current mainline implementation, $\Phi_{\theta}$ is typically instantiated as an LSTM encoder.

\subsection{Grouped Prediction Head}

Conditioned on $\mathbf{h}_k$, the prediction head produces future state increments:
\begin{equation}
\Delta \hat{\mathbf{Y}}_{k+1:k+H}
=
\Psi_{\theta}\!\left(\mathbf{h}_k\right)
\in \mathbb{R}^{H \times D_{\mathrm{out}}}.
\end{equation}

When enabled, the uncertainty head additionally predicts
\begin{equation}
\log \hat{\boldsymbol{\sigma}}^2_{k+1:k+H}
=
\Omega_{\theta}\!\left(\mathbf{h}_k\right)
\in \mathbb{R}^{H \times D_{\mathrm{out}}}.
\end{equation}

The grouped-head design reflects the current engineering prior that acceleration, angular rate, and velocity should remain distinguishable state groups during optimization and evaluation.

\subsection{Multi-Step Block Predictor}

For the current multi-step mainline, the model predicts a block of future increments with horizon $H>1$:
\begin{equation}
\left(
\Delta \hat{\mathbf{Y}}_{k+1:k+H},
\log \hat{\boldsymbol{\sigma}}^2_{k+1:k+H}
\right)
=
f_{\theta}\!\left(\mathbf{X}^{(L)}_k\right).
\end{equation}

Let the last proxy state inside the history window be denoted by $\hat{\mathbf{x}}_k$.
The predicted future state sequence is reconstructed by delta-cumsum rollout:
\begin{equation}
\hat{\mathbf{x}}_{k+j}
=
\hat{\mathbf{x}}_{k}
+
\sum_{i=1}^{j}\Delta \hat{\mathbf{x}}_{k+i},
\qquad j = 1,\dots,H.
\end{equation}

This contract is used by the current training and evaluation pipelines because it preserves direct comparability with existing long-rollout metrics and artifacts.

\subsection{Single-Step Transition Branch}

To move closer to an explicit state-transition formulation, the project also maintains a single-step branch with $H=1$:
\begin{equation}
\Delta \hat{\mathbf{x}}_{k+1}
=
f_{\theta}^{\mathrm{step}}\!\left(\mathbf{X}^{(L)}_k\right),
\qquad
\hat{\mathbf{x}}_{k+1}
=
\hat{\mathbf{x}}_{k} + \Delta \hat{\mathbf{x}}_{k+1}.
\end{equation}

Although the implementation still consumes a history window, the predictive semantics are much closer to
\begin{equation}
\hat{\mathbf{x}}_{k+1} = f_{\vartheta}\!\left(\hat{\mathbf{x}}_{k}, \mathbf{u}_{k}, \mathbf{c}_{k}\right),
\end{equation}
because the output is restricted to a one-step update.

\subsection{Optional Structured Modules}

The network may include additional structured modules:
\begin{itemize}
  \item \textbf{ThrusterLag}: transforms raw command history into effective control features,
  \item \textbf{HydroSSM}: provides latent hydrodynamic memory,
  \item \textbf{DampingHead}: injects damping-related inductive bias into predicted increments,
  \item \textbf{UncertaintyHead}: predicts diagonal log-variance terms.
\end{itemize}

Accordingly, one may interpret the overall model abstractly as
\begin{equation}
\left(
\Delta \hat{\mathbf{Y}},
\log \hat{\boldsymbol{\sigma}}^2
\right)
=
f_{\theta}\!\left(\mathbf{X}^{(L)}_k, \mathbf{r}_k\right),
\end{equation}
where $\mathbf{r}_k$ denotes optional structured latent features induced by these modules.

\subsection{Architectural Interpretation}

The current architecture should therefore be understood as a transition-model family with two active semantic modes:
\begin{itemize}
  \item a multi-step block predictor used to preserve evaluation continuity,
  \item a single-step transition branch used to approach solver-style deployment.
\end{itemize}

This distinction is important for both engineering and scientific writing:
the former is the current operational baseline, while the latter is the mathematically cleaner path toward controller and RL integration.
